% ############################################################################
\section{Dockerized Executables}

% Related documentation files:
% - helpers_root/docs/tools/docker/all.dockerized_flow.explanation.md
%   Core explanation of dockerized executable pattern
% - helpers_root/docs/tools/dev_system/all.docker_optimizer_container.how_to_guide.md
%   Specific example of dockerized executable for optimization
% - helpers_root/docs/tools/documentation_toolchain/all.latex_toolchain.how_to_guide.md
%   LaTeX as a dockerized executable
% - helpers_root/docs/tools/documentation_toolchain/all.render_images.explanation.md
%   Image rendering tools as dockerized executables
% - helpers_root/docs/tools/documentation_toolchain/all.notes_to_pdf.explanation.md
%   PDF generation as dockerized executable
% - helpers_root/docs/tools/all.llm_transform.how_to_guide.md
%   LLM-powered transformations as dockerized executables

% ============================================================================
\subsection{Motivation}

Software development teams frequently encounter the challenge of managing
diverse toolchains across heterogeneous environments. Installing and
maintaining developer utilities often leads to version conflicts, platform-specific
incompatibilities, and inconsistent behavior across development
machines, CI/CD pipelines, and production systems. These issues compound
when teams work across multiple operating systems (macOS, Linux, WSL) and
deployment contexts (local development, containerized environments,
cloud infrastructure).

Traditional approaches to tool management present several limitations.
Installing tools directly on host systems creates fragile dependencies
that differ across platforms. Packaging all utilities into a single monolithic
development container leads to bloated images, lengthy build times, and
dependency conflicts between tools. Version management becomes complex
when different projects require different tool versions, and onboarding
new team members involves extensive setup documentation that quickly becomes
outdated.

To address these challenges, we developed \textbf{Dockerized
Executables}---a pattern that packages each development tool in its own
slim, purpose-built container image. Instead of installing tools locally
or embedding them in a single development container, each utility runs
on-demand in an isolated environment with all dependencies pre-configured.
This approach provides clean separation between tools, enables
independent versioning and upgrades, reduces security surface area,
and ensures consistent behavior across all execution contexts.

% ============================================================================
\subsection{Summary}

Dockerized Executables solve a persistent problem in software
development: the challenge of maintaining consistent, reproducible
tooling across diverse environments. By packaging each tool in an isolated
container with a transparent wrapper, we eliminate installation
complexity, ensure consistent behavior, and maintain clean development
environments. While this approach requires Docker as a dependency, the
benefits in reproducibility, onboarding speed, and operational
simplicity justify this requirement. The pattern has proven effective across
diverse tooling needs, from document processing to diagram generation
to PDF compilation, and has become a foundational component of our development
infrastructure.

% ============================================================================
\subsection{Architecture and Design}

A Dockerized Executable consists of two primary components: a minimal
container image containing the tool and its dependencies, and a thin wrapper
script that handles execution logistics. From the developer's
perspective, invoking a Dockerized Executable is indistinguishable from
running a locally installed tool---the wrapper transparently manages containerization
details.

% ----------------------------------------------------------------------------
\subsubsection{Container image structure}

Each tool is packaged in a dedicated Docker image built on a minimal
base (e.g., \texttt{debian:stable-slim}). The image includes only the
dependencies required for that specific utility, keeping the image size
small and reducing attack surface. Images are version-pinned and
tagged to enable controlled upgrades and rollbacks. The container runs
with a non-root user by default and uses a standardized working directory
(\texttt{/src}) where the repository is mounted.

% ----------------------------------------------------------------------------
\subsubsection{Wrapper script responsibilities}

The wrapper script bridges between the user's environment and the containerized
tool. Key responsibilities include:

\begin{itemize}
  \item \textbf{Repository discovery}: Locates the repository root to establish
    a consistent mount point regardless of the user's current working directory.

  \item \textbf{Path translation}: Converts file paths between host/dev-container
    coordinates and container-internal paths (e.g., \texttt{docs/file.md}
    → \texttt{/src/docs/file.md}).

  \item \textbf{Container lifecycle management}: Pulls or builds the image
    if needed (leveraging Docker's caching), launches the container
    with appropriate mounts and permissions, and removes it after execution.

  \item \textbf{Permission preservation}: Runs the container with the caller's
    UID:GID to ensure generated files have correct ownership.

  \item \textbf{Environment isolation}: Whitelists and forwards only necessary
    environment variables, preventing unintended dependencies on host
    configuration.

  \item \textbf{Output handling}: Streams stdout/stderr in real-time and
    propagates the tool's exit code for proper CI/CD integration.
\end{itemize}

% ----------------------------------------------------------------------------
\subsubsection{Repository root as anchor point}

Using the repository root as the mount anchor is critical for
consistency. This design choice ensures:

\begin{itemize}
  \item Commands work identically regardless of the user's current subdirectory

  \item Paths remain stable across local development, dev containers, and
    CI environments

  \item The same command works in both children and sibling container patterns

  \item Minimal, predictable mounts prevent exposing unrelated host
    directories
\end{itemize}

% ============================================================================
\subsection{Execution Flow}

When a developer invokes a Dockerized Executable:

\begin{enumerate}
  \item The wrapper script is called with the tool's arguments

  \item The script discovers the repository root and normalizes input paths

  \item It ensures the tool's Docker image is available (pulling/building
    if necessary, leveraging cache otherwise)

  \item The repository is mounted into a new container at \texttt{/src}

  \item Paths are rewritten to container coordinates

  \item The tool executes inside the container with the provided arguments

  \item Outputs are written back to the mounted repository

  \item The container exits, and its exit code is propagated to the
    caller
\end{enumerate}

The entire process is transparent to the user, who sees only the tool's
output and return code.

% ============================================================================
\subsection{Container Execution Patterns}

Dockerized Executables must operate correctly whether invoked from the
host system or from within development containers. This requirement introduces
complexity in mounting and Docker daemon access.

% ----------------------------------------------------------------------------
\subsubsection{Children containers (Docker-in-Docker)}

In this pattern, a development container runs its own Docker daemon and
spawns tool containers as children. While flexible, this approach requires
elevated privileges (typically \texttt{--privileged}) and introduces
security concerns. It also adds overhead from running nested Docker daemons.

% ----------------------------------------------------------------------------
\subsubsection{Sibling containers (preferred)}

Our preferred approach mounts the host's Docker socket into the
development container, allowing it to launch sibling containers via the
host's Docker engine. The development container and tool container run
as peers, both managed by the host daemon. This pattern is more secure,
more efficient, and avoids privilege escalation. However, it requires
careful path management: mounts must reference host-visible paths, not
dev-container-internal paths. The wrapper script handles this translation
automatically.

% ============================================================================
\subsection{Practical Examples}

We have successfully Dockerized numerous development utilities:

\begin{itemize}
  \item \textbf{Document formatting}: Prettier for consistent code and
    Markdown formatting across the codebase

  \item \textbf{Document conversion}: Tools for converting between formats
    (e.g., DOCX to Markdown) without installing converters locally

  \item \textbf{Diagram rendering}: Mermaid, Graphviz, and TikZ renderers
    for generating diagrams from source specifications

  \item \textbf{LaTeX compilation}: Complete LaTeX toolchain for building
    PDFs without installing LaTeX distributions on development machines

  \item \textbf{LLM-powered transforms}: Utilities leveraging language
    models for structured content transformations
\end{itemize}

Each tool follows the same pattern: a single Python wrapper script, a minimal
Dockerfile, and standardized invocation semantics.

% ----------------------------------------------------------------------------
\subsubsection{Example invocation}

Converting a document requires only:

\begin{verbatim}
python dev_scripts_helpers/documentation/convert_docx_to_markdown.py \
-i docs/weekly_update.docx \
-o docs/weekly_update.md
\end{verbatim}

The wrapper handles image management, path translation, and execution
automatically. The result is identical to a local installation, but
with consistent behavior across all platforms.

% ============================================================================
\subsection{Benefits and Trade-offs}

The Dockerized Executable pattern provides several advantages:

\begin{itemize}
  \item \textbf{Reproducibility}: Identical tool behavior across all environments
    (macOS, Linux, WSL, CI)

  \item \textbf{Rapid onboarding}: Docker is the only prerequisite; no
    lengthy installation guides

  \item \textbf{Clean environments}: Development containers remain lean,
    containing only essential dependencies

  \item \textbf{Independent versioning}: Each tool can be upgraded or rolled
    back without affecting others

  \item \textbf{Reduced conflicts}: Isolated dependencies prevent version
    collisions between tools

  \item \textbf{Smaller security surface}: Purpose-built images contain
    only necessary components
\end{itemize}

However, this approach introduces some costs:

\begin{itemize}
  \item \textbf{Docker dependency}: The system requires Docker to be available
    and properly configured

  \item \textbf{Image storage}: Each tool requires disk space for its image,
    though images are typically small (tens to hundreds of MB)

  \item \textbf{Initial pull latency}: First-time execution requires pulling
    or building the image, though subsequent runs leverage Docker's
    cache

  \item \textbf{Slight overhead}: Container startup adds minimal latency
    (typically under 1 second) compared to native execution
\end{itemize}

For tools used frequently in development workflows, these trade-offs are
favorable. The reproducibility and consistency benefits far outweigh
the minimal overhead.

% ============================================================================
\subsection{Comparison with Language-Specific Solutions}

The Dockerized Executable pattern addresses a broader problem space than
language-specific dependency managers. Tools like \texttt{uv} (for
Python) provide fast, isolated dependency management within their
language ecosystem. They excel at version pinning and environment isolation
for language packages but are limited to that language's tooling.

Dockerized Executables complement such tools rather than replacing them.
They provide consistent execution for:

\begin{itemize}
  \item Tools written in languages not used in the main project

  \item System-level utilities with complex native dependencies

  \item Tools requiring specific OS packages or libraries

  \item Cross-language workflows requiring coordinated toolchains
\end{itemize}

Both approaches share the goal of reproducible, isolated tool
execution. The choice depends on whether the tool fits within a language
ecosystem's package manager or requires the broader isolation that containers
provide.

% ============================================================================
\subsection{Implementation Guidelines}

Building a Dockerized Executable follows a consistent pattern:

\begin{enumerate}
  \item \textbf{Create a minimal image}: Start with a small base image
    (\texttt{debian:stable-slim}, \texttt{alpine}, or language-specific
    slim variants). Install only necessary dependencies. Pin versions explicitly.
    Use a non-root user. Set \texttt{WORKDIR /src}.

  \item \textbf{Build the wrapper script}: Implement repository root discovery.
    Mount the repo at \texttt{/src}. Translate paths from host to
    container coordinates. Forward necessary arguments and environment
    variables. Run with caller's UID:GID. Stream logs and propagate exit
    codes.

  \item \textbf{Choose execution pattern}: Prefer sibling containers for
    security and efficiency. Handle path translation correctly for
    sibling pattern. Document any special requirements.

  \item \textbf{Test thoroughly}: Invoke from repository root and subdirectories.
    Test from host and dev containers. Verify path handling and file ownership.
    Assert on outputs and exit codes.

  \item \textbf{Document usage}: Provide clear examples. Document
    required environment variables if any. Note any platform-specific considerations.
\end{enumerate}

% ============================================================================
\subsection{Integration with Development Workflows}

Dockerized Executables integrate seamlessly into existing workflows:

\begin{itemize}
  \item \textbf{Local development}: Developers invoke tools via wrapper
    scripts, receiving immediate feedback

  \item \textbf{Pre-commit hooks}: Git hooks can invoke Dockerized Executables
    for formatting or validation

  \item \textbf{CI/CD pipelines}: GitHub Actions and other CI systems invoke
    the same wrapper scripts, ensuring consistency

  \item \textbf{Automated workflows}: Tools can be orchestrated via make,
    invoke, or other task runners

  \item \textbf{Interactive use}: Developers can use tools ad-hoc during
    development without environment concerns
\end{itemize}

This consistency across contexts eliminates "works on my machine" issues
and simplifies troubleshooting.

% ############################################################################
\section{Docker container release flow}

% Related documentation files:
% - docs/infra/ck.development_stages.explanation.md
%   Explains dev/prod/local container stages
% - docs/infra/ck.elastic_container_service.explanation.md
%   ECS setup for container deployment
% - docs/infra/ck.ecs_flows.reference.md
%   Reference for ECS container flows
% - docs/infra/ck.aws_ecr.explanation.md
%   AWS ECR for container registry management
% - docs/infra/ck.airflow_dags_release.how_to_guide.md
%   Airflow DAG release process coordinated with containers
% - helpers_root/docs/tools/dev_system/all.devops_docker_auto_release.explanation.md
%   Automated Docker container release workflows
% - helpers_root/docs/tools/dev_system/all.devops_docker.how_to_guide.md
%   Manual Docker release procedures
% - helpers_root/docs/tools/dev_system/all.devops_docker.reference.md
%   Complete reference for container release system
% - docs/build/ck.gh_workflows.explanation.md
%   GitHub workflows for building containers
% - docs/build/ck.gh_workflows.reference.md
%   Reference for all build workflows
% - docs/bwd/ck.bwd_app_dev_image_release.how_to_guide.md
%   Development image release process
% - docs/bwd/ck.bwd_app_backend_release.how_to_guide.md
%   Backend production container release
% - docs/trading_ops/ck.release.explanation.md
%   Release process including container updates
% - docs/trading_ops/ck.release.how_to_guide.md
%   Step-by-step release guide
% - docs/trading_ops/ck.dags_release.how_to_guide.md
%   DAG release synchronized with container versions
% - docs/all.dev_system_release.how_to_guide.md
%   Overall development system release process

% ============================================================================
\subsection{Motivation and Architecture}
Modern software systems require consistent, reproducible deployment environments
across development, staging, and production. Traditional deployment
approaches struggle with environmental drift, dependency conflicts,
and versioning complexity. Docker containers provide isolation and reproducibility,
but managing container lifecycles—building, testing, versioning, and deploying—introduces
operational overhead. Our Docker container release flow addresses
these challenges through automated workflows that enforce quality
gates, maintain version consistency, and provide staged rollout
capabilities. The system balances developer velocity with operational safety,
enabling rapid iteration while preventing production incidents. The architecture
distinguishes between multiple container types and deployment stages:
\begin{itemize}
  \item \textbf{Development images} (\texttt{dev}): Used for day-to-day
    development, containing all development dependencies, debugging
    tools, and test frameworks. Updated frequently (daily or on-demand).

  \item \textbf{Production images} (\texttt{prod}): Optimized for
    deployment, containing only runtime dependencies with minimal attack
    surface. Updated on a controlled schedule (bi-weekly or when critical
    fixes are required).

  \item \textbf{Local images} (\texttt{local}): Developer-specific
    builds for testing container changes before promoting to dev. Never
    pushed to registries.
\end{itemize}
Each image stage has corresponding task definitions in AWS ECS (Elastic
Container Service) for orchestrated deployments. The entire flow is managed
through GitHub Actions workflows that run on push, schedule, or manual
dispatch.
% ============================================================================
\subsection{Container Build Pipeline}
The container build pipeline transforms source code and dependencies
into versioned, tested Docker images ready for deployment. The process
is automated through GitHub Actions workflows that run in response to
code changes or manual triggers.
% ----------------------------------------------------------------------------
\subsubsection{Development Image Workflow}
Development images are built and released through the \texttt{build\_image.dev.yml}
workflow. This workflow:
\begin{enumerate}
  \item Triggers on:
    \begin{itemize}
      \item Push to master branch

      \item Manual workflow dispatch with optional version override

      \item Scheduled runs (daily for routine updates)
    \end{itemize}

  \item Checks out the repository at the specified commit or branch

  \item Reads the current version from \texttt{changelog.txt}, which maintains
    a human-readable history of container changes with semantic versioning
    (major.minor.patch)

  \item Builds the image using the Dockerfile in the repository root,
    installing development dependencies, test frameworks, and
    debugging tools

  \item Tags the image with:
    \begin{itemize}
      \item Version tag (e.g., \texttt{1.2.3})

      \item \texttt{dev} tag pointing to the latest development
        version

      \item Commit SHA for precise traceability
    \end{itemize}

  \item Pushes the image to Amazon ECR (Elastic Container Registry)

  \item Updates environment variables and configuration to reference the
    new image version

  \item Optionally triggers downstream workflows to update dependent systems
\end{enumerate}
The development image build is relatively fast (typically 10-20 minutes)
because it leverages Docker layer caching and does not require
extensive optimization passes.
% ----------------------------------------------------------------------------
\subsubsection{Production Image Workflow}
Production images follow a similar but more rigorous process through \texttt{build\_image.cmamp.yml}
(or equivalent production workflow). Key differences include:
\begin{itemize}
  \item \textbf{Stricter trigger conditions}: Only builds on manual dispatch
    or scheduled releases, never on every push. This prevents immature
    code from reaching production workflows.

  \item \textbf{Multi-stage builds}: Uses Docker multi-stage builds to
    minimize final image size. Build dependencies are discarded, and only
    runtime requirements remain in the final image.

  \item \textbf{Security scanning}: Integrates vulnerability scanning (e.g.,
    Trivy, AWS ECR scanning) to detect known CVEs in base images and
    dependencies before release.

  \item \textbf{Optimization passes}: Removes unnecessary files,
    compresses layers, and ensures minimal attack surface

  \item \textbf{Comprehensive testing}: Runs full test suites (fast, slow,
    and sometimes superslow) to validate the production build before release

  \item \textbf{Version locking}: Production images are immutable once
    tagged. Rollbacks occur by redeploying a previous version, not by modifying
    existing images.
\end{itemize}
Production image builds are slower (30-60 minutes) due to these
additional steps, but the investment ensures reliability and security.
% ----------------------------------------------------------------------------
\subsubsection{Version Management}
Container versions are managed through \texttt{changelog.txt} files that
serve as the single source of truth for the current version. Each
entry includes:
\begin{itemize}
  \item Version number (semantic versioning: major.minor.patch)

  \item Release date

  \item Summary of changes (features, fixes, dependency updates)

  \item Author information
\end{itemize}
When a developer needs to release a new version, they:
\begin{enumerate}
  \item Update \texttt{changelog.txt} with the new version and
    description

  \item Commit the change to the repository

  \item Trigger the build workflow (automatically on push to master or
    manually via workflow dispatch)

  \item The workflow reads the version from \texttt{changelog.txt} and
    tags the image accordingly
\end{enumerate}
This approach ensures version changes are tracked in version control and
reviewed during code review, preventing accidental or undocumented
version bumps.
% ============================================================================
\subsection{Task Definition Management with ECS}
AWS Elastic Container Service (ECS) orchestrates containerized applications
using task definitions that specify which Docker images to run,
resource allocations, networking configuration, and environment
variables. Our system manages separate task definitions for preproduction
and production environments, enabling staged rollouts and safe testing
of changes.

% ----------------------------------------------------------------------------
\subsubsection{Preproduction Task Definition Release}
The \texttt{release\_new\_ecs\_preprod\_task\_definition.yml} workflow
updates the preproduction ECS task definition. This workflow:
\begin{enumerate}
  \item Triggers on manual dispatch, allowing controlled releases
    after image builds complete

  \item Fetches the current task definition from AWS ECS

  \item Updates the image reference to point to the newly built development
    image (specified by version or tag)

  \item Modifies environment variables if needed (e.g., database
    connection strings for preproduction)

  \item Registers the new task definition revision with ECS

  \item Optionally forces a new deployment of running services using this
    task definition, causing ECS to restart tasks with the new image

  \item Logs the task definition ARN and revision number for traceability
\end{enumerate}
Preproduction task definitions use development images and connect to
isolated resources (databases, message queues) to avoid impacting production
systems. This environment serves as the final validation stage before
production deployment.

% ----------------------------------------------------------------------------
\subsubsection{Production Task Definition Release}
The \texttt{release\_new\_ecs\_prod\_task\_definition.yml} workflow
follows an identical structure but with additional safeguards:
\begin{itemize}
  \item Requires explicit manual approval (via GitHub environment
    protection rules) before proceeding

  \item Only references production-tagged images that have passed all quality
    gates

  \item Logs all changes to an audit trail for compliance and debugging

  \item Supports blue-green deployments and canary releases through ECS
    service configuration

  \item Monitors deployment progress and can automatically roll back if
    health checks fail
\end{itemize}
Production releases typically occur bi-weekly on a fixed schedule unless
critical patches are required. This cadence balances feature velocity
with operational stability.

% ============================================================================
\subsection{Airflow DAG Release Process}
Airflow DAGs (Directed Acyclic Graphs) define workflows for data processing,
model training, and automated trading operations. DAG code changes
often require updated Docker images containing new dependencies or
business logic. The DAG release process coordinates code changes with container
updates.

% ----------------------------------------------------------------------------
\subsubsection{Release Workflow for Trading DAGs}
The release procedure for Airflow DAGs follows a structured, multi-step
process:
\begin{enumerate}
  \item \textbf{Create and test a DAG}: When updating DAG logic,
    create a test DAG with a task-specific name (e.g., \texttt{test.tokyo.shadow\_trading.C11a.config1\_CmTask8080}).
    This naming convention embeds the task number for traceability.

  \item \textbf{Validate with preproduction task definition}: Execute the
    test DAG using the existing preproduction ECS task definition. Do
    not update the image yet—this isolates DAG logic changes from
    environment changes. Run for a minimum duration (e.g., 30 minutes)
    to ensure stability.

  \item \textbf{Verify all tasks succeed}: Check that every task in the
    DAG completes successfully without errors or warnings. Review logs
    for unexpected behavior.

  \item \textbf{Propagate changes to production DAGs}: Once validated,
    apply the changes to all relevant production DAGs in \texttt{datapull/airflow/dags/trading/shadow\_trading/}.

  \item \textbf{Update shared utilities}: If new utility functions were
    added to \texttt{datapull/airflow/dags/airflow\_utils/trade\_exec/utils.py},
    copy them to the shared Airflow utilities directory (\texttt{/data/shared/airflow/dags/airflow\_utils/trade\_exec/utils.py})
    after validation.

  \item \textbf{Merge and document}: Open a pull request for review.
    Once approved, merge to master. Create a release issue titled ``Release
    shadow trading DAGs — \{date\}'' listing affected DAGs and links
    to merged PRs.

  \item \textbf{Image update if needed}: If the DAG changes require a new
    Docker image (e.g., new Python packages, updated business logic), note
    this in the release issue and coordinate with the production release
    schedule.
\end{enumerate}
This process ensures DAG changes are validated independently before
deployment and that dependencies between DAGs, images, and
infrastructure are clearly managed.

% ----------------------------------------------------------------------------
\subsubsection{Release Schedule}
Shadow trading DAGs follow a bi-weekly release cadence when image updates
are required. When only DAG logic changes (no new dependencies), releases
occur weekly. This schedule provides predictability for operations
teams while allowing rapid iteration on trading strategies.

% ============================================================================
\subsection{Feature Release Communication and Documentation}
Releasing new features or significant system changes requires more than
code deployment. Effective communication ensures users understand what
changed, why it matters, and how to use new capabilities. Our process integrates
documentation updates and stakeholder communication into the release workflow.

% ----------------------------------------------------------------------------
\subsubsection{Pre-Release Internal Testing}
Before broadly announcing a feature, we conduct internal testing with a
small group:
\begin{enumerate}
  \item \textbf{Draft communication}: Write a detailed description of the
    feature in a Google Doc, explaining what changed, why, expected
    user actions, and links to documentation or dashboards.

  \item \textbf{Team review}: Share the draft with the feature development
    team for feedback on technical accuracy and clarity.

  \item \textbf{Pilot rollout}: Send the communication to a limited group
    (feature team or power users) and ask for feedback on both the feature
    and the communication.

  \item \textbf{Iterate}: Incorporate feedback and resolve any issues discovered
    during the pilot.
\end{enumerate}
This approach catches problems early when they're easier to fix and
ensures communication is clear before reaching a wider audience.
% ----------------------------------------------------------------------------
\subsubsection{Updating Central Release Documentation}
After internal validation, update the master release tracking document
(\textit{Dev System Features} Google Doc) with:
\begin{itemize}
  \item Feature title

  \item Release date

  \item Summary of changes

  \item Links to detailed documentation, dashboards, or related issues
\end{itemize}
This central document serves as a historical record of system
evolution and a reference for onboarding new team members.
% ----------------------------------------------------------------------------
\subsubsection{Communication Template}
Release announcements follow a structured template:
\begin{verbatim}
Subject: [Feature Name] Now Available
Hi Team,
We've integrated [feature name] to [solve what problem].
What's Covered?
- Key capability 1
- Key capability 2
- Integration points and usage
Current State:
- Metrics or dashboards (with access links)
- Example visuals or screenshots
Next Steps:
- Recommended actions for users
- How to get help or provide feedback
For detailed documentation, see: [link]
Thanks,
[Team]
\end{verbatim}
This format ensures consistency, provides actionable information, and
reduces confusion about new features.
% ============================================================================
\subsection{Integration with CI/CD and Quality Gates}
The Docker release flow integrates tightly with our continuous integration
and deployment (CI/CD) pipelines. Quality gates at each stage prevent
broken code from progressing through the release pipeline.
% ----------------------------------------------------------------------------
\subsubsection{Automated Testing in Workflows}
Build workflows incorporate testing at multiple points:
\begin{itemize}
  \item \textbf{Pre-build validation}: Lint checks, code formatting, and
    static analysis run before building the image. If these fail, the build
    aborts immediately.

  \item \textbf{Post-build testing}: After the image is built, fast
    tests run inside the new container to validate that code executes
    correctly in the containerized environment. This catches environment-specific
    issues (e.g., missing dependencies, incorrect paths).

  \item \textbf{Integration testing}: For production builds,
    integration tests verify interactions between components (databases,
    APIs, message queues) using the new image.

  \item \textbf{Security scanning}: Vulnerability scanners analyze the
    image for known CVEs. High-severity vulnerabilities block the
    release until patched.
\end{itemize}
Each gate can block progression, ensuring only validated images reach deployment
stages.
% ----------------------------------------------------------------------------
\subsubsection{Rollback Capabilities}
Despite rigorous testing, issues sometimes arise in production. The
system supports rapid rollback:
\begin{itemize}
  \item \textbf{Immutable image tags}: Each version is permanently tagged
    in ECR. Rollback is as simple as updating the task definition to
    reference a previous version.

  \item \textbf{Git-tracked configuration}: Task definition changes are
    committed to the repository. To roll back, revert the commit and re-run
    the release workflow.

  \item \textbf{Automated health checks}: ECS monitors container health.
    If new tasks fail health checks repeatedly, ECS can automatically stop
    the rollout and maintain the previous version.

  \item \textbf{Blue-green deployments}: For critical services, we
    maintain parallel environments (blue and green). New versions deploy
    to the inactive environment, and traffic switches only after
    validation succeeds.
\end{itemize}
These mechanisms minimize downtime and reduce the blast radius of problematic
releases.
% ============================================================================
\subsection{Buildmeister Dashboard and Monitoring}
The Buildmeister Dashboard provides real-time visibility into the
state of all GitHub Actions workflows across all repositories. This centralized
view enables rapid identification of failures and simplifies daily operational
oversight. The dashboard displays:
\begin{itemize}
  \item \textbf{Workflow status}: Current state of all workflows (success,
    failed, in progress) across repositories

  \item \textbf{Links to runs}: Direct links to failing workflow runs for
    immediate investigation

  \item \textbf{Test reports}: Links to Allure reports showing test history,
    flaky tests, and failure trends

  \item \textbf{Coverage reports}: Latest code coverage metrics with drill-down
    by module

  \item \textbf{Pull request counts}: Number of open PRs per repository,
    indicating review backlog
\end{itemize}
The dashboard updates hourly and is published to an S3-hosted static site
accessible to all team members. Historical versions are retained for
analysis of workflow stability over time. This tool is central to the Buildmeister
role, which rotates weekly among team members. The Buildmeister monitors
the dashboard, triages failures, and ensures builds remain green.
% ============================================================================
\subsection{Summary}
Our Docker container release flow automates the journey from source
code to deployed containers while enforcing quality gates, maintaining
version consistency, and enabling safe rollouts. Separate workflows for
development and production images balance velocity with stability.
Integration with ECS task definitions provides orchestrated
deployments with rollback capabilities. The Airflow DAG release process
coordinates workflow updates with container changes. Structured
communication and documentation ensure stakeholders understand changes.
Finally, the Buildmeister Dashboard and automated monitoring provide
visibility into system health and enable rapid response to issues. This
system transforms container management from a manual, error-prone
process into a reliable, auditable pipeline. Developers can iterate
rapidly with development images while operations teams maintain
production stability through controlled releases. The architecture is designed
for growth: adding new services or repositories requires minimal
configuration, as all core workflows are reusable and parameterized.
Quality gates prevent broken code from reaching production, and comprehensive
monitoring ensures issues are detected and resolved quickly.
